% Fragmento conservado; véase MANIFIESTO_ANTECEDENTES_LECTURA.json.
\subsection{Reversibilidad de la constante racional y cambio de fase}

La constante racional \(\kappa_{\mathrm{per}}\) determina íntegramente
el registro \(K\) cuando se conservan la base, la longitud y el origen
de lectura. Su función numérica y su función estructural se relacionan
mediante una inversión exacta.

\begin{proposition}[Recuperación integral y transformación cíclica]
Sea \(I(K)=\sum_{j=1}^{12}K_j1000^{12-j}\). Entonces
\[
 I(K)=(1000^{12}-1)\kappa_{\mathrm{per}}.
\]
La expansión entera de \(I(K)\) en doce bloques de base \(1000\)
recupera todos los \(K_j\), incluidos los ceros iniciales de cada bloque.
Si \(\rho(K)=(K_2,\ldots,K_{12},K_1)\), se cumple
\begin{equation}
 \kappa(\rho K)=1000\kappa(K)-K_1.
 \label{eq:k-rotacion-racional}
\end{equation}
\end{proposition}
\begin{proof}
La primera identidad invierte la definición de la evaluación periódica.
La división euclídea iterada por \(1000\), con doce posiciones fijadas,
recupera el vector. Por otra parte,
\[
 I(\rho K)=1000I(K)-K_1(1000^{12}-1).
\]
Dividir por \(1000^{12}-1\) da \eqref{eq:k-rotacion-racional}.
\end{proof}

La recuperación compone, por tanto, \(\kappa\mapsto K\mapsto U\)
y las transformaciones \(D_3K,D_4K,Q\) ya construidas. La evaluación
racional conserva el registro firmado y las diferencias cíclicas que
dependen de él. La profundidad total y los operadores de transporte de
una historia permanecen en sus coordenadas propias: la inversión anterior
tiene por dominio el registro integral dodecafásico.

Escribamos \(\kappa_j=\kappa(\rho^jK)\), con índices módulo \(12\).
La misma identidad da
\[
 K_{j+1}=1000\kappa_j-\kappa_{j+1},\qquad
 \sum_{j=0}^{11}\kappa_j=\frac{\sum_{j=1}^{12}K_j}{999}
 =\frac{6263}{999}.
\]
La fase actúa de manera exacta sobre la constante: su valor es fijo
en la lectura canónica, mientras sus rotaciones siguen una ley afín.
La suma de la órbita es invariante bajo la elección del origen.

También la incidencia posterior dispone de una expresión en estas
coordenadas. El umbral cúbico determina
\[
 B_K=\{j+1:1000\kappa_j-\kappa_{j+1}\ge729\}
 =\{4,6,9,10\}.
\]
Se recupera la tétrada a partir de la órbita racional porque ésta
recupera primero las componentes enteras. El soporte firmado asociado
a \(\alpha\) utiliza además sus registros regionales y la normalización
de frontera. La relación entre coordenada e incidencia queda así
expresada mediante operaciones identificables.

Esta es una definición operacional de la función de \(K\): proporciona
una coordenatización integral reversible del registro orientado, cuya
evaluación periódica y cuyas funciones de incidencia participan en
construcciones posteriores. El valor, la transformación y el dominio
se conservan conjuntamente.

